
\documentclass[interno]{gpt}

\titolo{Analisi dei capitolati}
\data{2026-10-15}
\orario{17:30}{18:00}
\luogo{Discord}
\redattore{mazzurana}
\verificatore{corro}
\componenti{biasini,cescobolla,corro,guglielmo,mazzurana,spagnoli}

\modifica{1.0.0}{2026-10-16}{mazzurana}{corro}{Verifica e approvazione}
\modifica{0.1.0}{2026-10-15}{mazzurana}{corro}{Prima stesura}

\begin{document}

\section{Ordine del giorno}
\begin{enumerate}
  \item Definizione dei prossimi passi e pianificazione delle attività.
  \item Ripartizione della stesura della documentazione in corso:
    \begin{itemize}
      \item Verbale esterno dell'incontro odierno;
      \item Documento di analisi dei capitolati;
      \item Preventivo costi e impegni orari;
      \item Lettera di presentazione collettiva per l'incontro straordinario.
    \end{itemize}
\end{enumerate}
\clearpage
\section{Svolgimento}
\subsection{Definizione dei prossimi passi}
Il gruppo si è riunito per coordinare le attività a seguito delle ultime interazioni con i proponenti e definire i documenti prioritari da redigere e revisionare nei prossimi giorni. Sono stati individuati quattro task principali legati alla formalizzazione delle decisioni e alla preparazione per l'incontro straordinario.

\subsection{Organizzazione e ripartizione dei compiti}
Per garantire il rispetto delle scadenze e una distribuzione equilibrata del carico di lavoro, i compiti sono stati così suddivisi:
\begin{itemize}
  \item \textbf{Verbale esterno:} documentazione formale dell'incontro svolto in data odierna;
  \item \textbf{Analisi dei capitolati:} consolidamento delle valutazioni comparative e delle risposte ricevute dai proponenti;
  \item \textbf{Preventivo:} stima dei costi, delle risorse e dell'impegno orario complessivo e individuale;
  \item \textbf{Lettera di presentazione:} redazione congiunta del testo formale da presentare in vista dell'incontro straordinario.
\end{itemize}

\section{Decisioni}
\begin{decisioni}
  \decisione{Procedere in parallelo con la redazione della documentazione assegnando ciascun documento a coppie dedicate.}
  \decisione{La lettera di presentazione per l'incontro straordinario sarà redatta in forma collettiva con il contributo e l'approvazione dell'intero gruppo.}
\end{decisioni}

\section{Azioni da intraprendere}
\begin{azioni}
  \azione[\#1]{Redazione del verbale esterno dell'incontro odierno.}{mazzurana, corro}{28/10/2026}
  \azione[\#2]{Redazione del documento di analisi dei capitolati.}{guglielmo, biasini}{29/10/2026}
  \azione[\#3]{Stesura del preventivo (costi e pianificazione oraria).}{spagnoli, cescobolla}{29/10/2026}
  \azione[\#4]{Stesura della lettera di presentazione per l'incontro straordinario.}{Tutti}{30/10/2026}
\end{azioni}

\section{Prossima riunione}
Giovedì 29 ottobre 2026 alle 17:30 su Discord, per la verifica congiunta dei documenti prima della consegna.

\end{document}
